\documentclass[journal]{IEEEtran}
\usepackage{amsmath,amssymb,amsthm}
\usepackage{booktabs}
\usepackage{graphicx}
\usepackage{url}
\usepackage[hidelinks]{hyperref}

\newtheorem{definition}{Definition}
\newtheorem{lemma}{Lemma}
\newtheorem{theorem}{Theorem}
\newcommand{\Rq}{\mathcal{R}_q}
\newcommand{\norm}[1]{\lVert #1 \rVert}

\begin{document}

\title{In-Band Lattice Zero-Knowledge Proofs for Compressed-Domain\\ H.264 Video Steganography}

\author{Anonymous authors (double-blind review)}

\maketitle

\begin{abstract}
Proof-carrying steganography hides a payload in a cover medium together with a
zero-knowledge proof about it, so that a verifier can check the claim from the
medium alone. Existing systems either rely on pairing-based SNARKs, which do
not survive quantum adversaries, or keep the proof outside the medium. We
present a system that embeds a post-quantum zero-knowledge proof directly in
the CAVLC residuals of an H.264 Baseline bitstream and verifies it from the
received video, without side files, SEI messages, or the cover. The proof is a
Fiat--Shamir-with-aborts proof of knowledge of the opening of a BDLOP lattice
commitment, instantiated at an estimated 138--140-bit classical (125--127-bit
quantum) core-SVP level. Its statement binds a carrier-normalized digest of
the video, the verifier-derived carrier positions, a one-use verifier
challenge and an issuer-signed relation registry; carrier normalization breaks
the circularity between a proof and the video that contains it. Two modes are
supported: a 7\,589-byte envelope proving possession of a registered
credential, and a 14\,309-byte envelope that commits to a hidden payload which
can later be disclosed or linked to another video without revealing it. We
report capacity on standard sequences, end-to-end embedding and video-only
verification, decoded quality, and classical and learned steganalysis.
\end{abstract}

\begin{IEEEkeywords}
Video steganography, H.264/AVC, CAVLC, zero-knowledge proofs, lattice-based
cryptography, post-quantum security, media provenance.
\end{IEEEkeywords}

\section{Introduction}

Compressed-domain steganography modifies quantized transform coefficients of
an encoded video, avoiding a decode/re-encode cycle and keeping the stream
standard-compliant~\cite{ref:h264stego}. Recent work attaches cryptographic
evidence to hidden data: a zero-knowledge proof that the hidden payload
satisfies a relation, carried with or inside the medium~\cite{ref:zkwagon}.
Two problems remain open for video.

First, deployed proof systems for this purpose are pairing-based. Groth16
proofs are 128--192 bytes and easy to hide, but their soundness and the
security of the commitments they prove about fall to Shor's algorithm. Media
provenance records are expected to outlive that horizon.

Second, a proof that lives in a side file or a metadata box is trivially
stripped or swapped, and a proof inside the video faces a circularity: the
proof cannot contain a hash of the bitstream that contains the proof.

We address both. Our contributions are:
\begin{itemize}
\item an in-band design in which a lattice zero-knowledge proof travels inside
  H.264 CAVLC residuals and is verified from the stego video alone, with a
  carrier-normalized video commitment that removes the proof/video
  circularity (Section~\ref{sec:construction});
\item a statement that binds the verifier-recomputed video commitment, carrier
  positions, a one-use verifier challenge and an ML-DSA-signed relation
  registry, giving replay protection and verifier-side trust anchoring;
\item two proof modes (registered credential and in-band hidden payload),
  plus linkage proofs showing that two videos, or a video and a credential,
  hide the same payload;
\item a concrete parameter set with security estimates, known-answer vectors,
  and an evaluation of capacity, quality and steganalysis on standard
  sequences (Section~\ref{sec:evaluation}).
\end{itemize}

\section{Related work}
\label{sec:related}

\paragraph{Proof-carrying media} PhotoProof~\cite{ref:photoproof} introduced
proof-carrying data for image authentication; ZK-IMG, VerITAS and
VIMz~\cite{ref:zkimg,ref:veritas,ref:vimz} prove image transformations with
SNARKs or folding schemes; Eva~\cite{ref:eva} proves the authenticity of H.264
video edits with incrementally verifiable computation. All keep the proof
outside the medium and use pre-quantum assumptions. ZK-WAGON~\cite{ref:zkwagon}
hides a zk-SNARK of model provenance inside an image by spatial LSB
embedding; it is the closest in-band precedent but is image-domain and not
post-quantum. Zero-knowledge watermark detection~\cite{ref:adelsbach} proves
the presence of a watermark without revealing it.

\paragraph{Lattice zero-knowledge} BDLOP commitments~\cite{ref:bdlop} with
Fiat--Shamir-with-aborts openings~\cite{ref:lyu09,ref:lyu12} underlie the
exact proofs of Attema et al.~\cite{ref:als20}, Esgin et al.~\cite{ref:ens20}
and Lyubashevsky et al.~\cite{ref:lnp22}. General lattice SNARKs such as
LaBRADOR, LatticeFold, RoKoko or SALSAA target succinctness and, in their
published form, do not provide zero knowledge; designated-verifier lattice
SNARKs (ISW21) rely on a knowledge assumption that is invalid against quantum
attackers. We surveyed these as backends and report why none fit
(Section~\ref{sec:discussion}).

\section{Background}

\paragraph{H.264 CAVLC} In Baseline profile, residual $4\times4$ blocks are
entropy-coded with CAVLC. The carrier rule used here selects, in each
eligible luma block, the first AC coefficient with $|\mathit{level}|\ge 4$ that
is not a trailing one, and writes one bit into its magnitude parity by a
change of $\pm1$. The change preserves the eligibility of the coefficient, so
a blind extractor reapplies the same rule; a patchability check rejects
blocks whose re-encoding would change the CAVLC syntax.

\paragraph{BDLOP commitments} Let $\Rq=\mathbb{Z}_q[X]/(X^N+1)$. With public
$B_1=[I_n\,|\,B_1']$ and $b_2=[0^n\,|\,1\,|\,b_2']$, a commitment to
$m\in\Rq$ with ternary randomness $r$ is $t_1=B_1r$, $t_2=\langle b_2,r\rangle+m$.
It is hiding under Module-LWE and binding under Module-SIS~\cite{ref:bdlop}.

\section{System and threat model}

Four roles interact: an \emph{issuer} signs a relation registry (ML-DSA-65);
a \emph{holder} produces the stego video; a \emph{verifier} issues one-use
session challenges and verifies; an \emph{adversary} observes, replays, edits
and re-embeds videos and may run steganalysis. The verifier trusts only its
pinned issuer key, parameter set and carrier profile, and its own session
store. It trusts nothing read from the video before verification succeeds.

We aim for (i) \emph{soundness}: an accepted video was produced by a party
knowing an opening of the bound commitment, for this video and session;
(ii) \emph{zero knowledge}: the video reveals nothing about the payload or the
opening beyond the statement; (iii) \emph{integrity}: any change outside the
carriers, and any change to the carriers, is rejected; (iv) \emph{replay
resistance}; and (v) \emph{undetectability} against passive steganalysis.
Camera provenance and codec correctness are out of scope.

\section{Construction}
\label{sec:construction}

\subsection{Proof system}
The prover samples $y$ uniformly in $[-\gamma+1,\gamma]^{kN}$, computes
$w=B_1y$, derives $c=\mathsf{SampleInBall}(H(\mathit{pid},x,t,w))$ with $\kappa$
coefficients $\pm1$, and outputs $(H(\cdot),z=y+cr)$ if
$\norm{z}_\infty\le B=\gamma-\kappa-1$, retrying otherwise. The verifier
recomputes $w'=B_1z-ct_1$ and checks the hash. The full algorithm, encodings
and domain tags are specified in the accompanying specification.

\subsection{Statement and carrier-normalized commitment}
The statement $x$ is a fixed-width encoding of the relation identifier,
parameter identifier, registry root and epoch, session challenge, codec
policy, carrier profile, carrier count, carrier-position hash and video
commitment $V$. $V$ hashes the bitstream with every carrier coefficient
replaced by a canonical value. Carrier positions are fixed from a placeholder
of the final envelope length before the proof is computed, so $V$ does not
depend on the proof bits. The verifier recomputes every field of $x$ from the
received video and its own configuration.

\subsection{Envelope, registry and sessions}
The envelope is
$\texttt{"ZKLB"}\|v\|\mathit{pid}\|N_s\|t\|\pi$ (in-band) or
$\texttt{"ZKLR"}\|v\|\mathit{pid}\|N_s\|\pi$ (registered), where in the
registered mode $t$ is a credential whose hash is part of the registry entry
signed by the issuer. It is split into fixed-size chunks with CRC framing and
embedded segment by segment. The session challenge $N_s$ is consumed
atomically after successful verification.

\subsection{Linkage}
An equality proof (the linear-relation proof of BDLOP) shows that two
commitments hide the same message; it costs two opening proofs and binds both
commitments and both statements.

\section{Security analysis}
\label{sec:security}

\begin{theorem}[informal]
In the random-oracle model, the proof is complete with per-attempt acceptance
$((2B+1)/2\gamma)^{kN}$, honest-verifier zero knowledge for accepted
transcripts, and special sound with extractor output a relaxed opening
$(\bar r,f)$ with $B_1\bar r=ft_1$, $\norm{\bar r}_\infty\le2B$ and $f$
invertible. Two relaxed openings to different messages yield a Module-SIS
solution of $\ell_\infty$ norm at most $8\kappa B$; hiding reduces to
Module-LWE of rank $\lambda$ with $n+1$ samples.
\end{theorem}

The proofs follow the published arguments for ML-DSA and BDLOP with uniform
masking; the gap in the Fiat--Shamir-with-aborts reprogramming argument is
treated by Barbosa et al.~\cite{ref:barbosa} and Devevey et al.~\cite{ref:devevey}.
We call the scheme plausibly post-quantum; we do not claim a tight
quantum-random-oracle proof.

\begin{table}[t]
\centering
\caption{Parameter set \texttt{bdlop-open-128-v2} and core-SVP estimates
(CRYSTALS security-estimates, commit \texttt{75c26949}; M-LWE dual hybrid
from the lattice-estimator, commit \texttt{e4333d9}).}
\label{tab:params}
\small\setlength{\tabcolsep}{4pt}
\begin{tabular}{ll}
\toprule
$N$, $q$ & $256$, $2^{30}-35$ ($q\equiv5 \bmod 8$) \\
$n$, $\lambda$, $k$ & $6$, $7$, $14$ \\
$\kappa$, $\gamma$, $B$ & $60$, $2^{17}$, $2^{17}-61$ \\
Challenge space & $2^{257}$ \\
Commitment / opening proof & 6\,720 B / 8\,096 B \\
Acceptance per attempt & 0.191 \\
Equality (linkage) proof & 16\,160 B \\
Max.\ in-band payload & 766 B \\
M-SIS ($\ell_\infty\le8\kappa B\approx2^{25.9}$) & $\beta=480$: 140 / 127 bits \\
M-LWE (rank $\lambda$, $n{+}1$ samples) & dual hybrid $\beta=564$ / $567$: 164 / 150 bits \\
\quad (primal uSVP) & $\beta=580$: 169 / 153 bits \\
\multicolumn{2}{l}{\footnotesize $\beta$: BKZ block size; core-SVP, classical / quantum.} \\
\bottomrule
\end{tabular}
\end{table}


\section{Evaluation}
\label{sec:evaluation}

\subsection{Setup}
All measurements were taken on a Windows 11 host with an Intel i7-12700H
(14 cores, 16\,GB RAM). Sequences are standard Xiph test clips (CIF
$352\times288$, 300 frames; 4CIF $704\times576$, 600 frames), encoded with x264
in Baseline/CAVLC, all-intra, at a fixed QP. The proof system and the H.264
analysis are pure Python. Unless stated otherwise the registered mode is used
(7\,589-byte envelope) and the envelope is split into 100-frame segments with a
chunk size chosen from the measured capacity of each clip. One run in each
mode also replays the consumed session and verifies a copy with one flipped
bit; both are rejected.

\subsection{Capacity}
Table~\ref{tab:capacity} lists the stable, patchability-checked carriers per
100 frames. At QP\,22 CIF clips offer 10--51\,kbit per 100 frames: textured
content (Container, Deadline) carries the registered envelope in 200 frames,
flat content (Akiyo, Foreman) cannot carry it in 300 frames. Capacity falls
steeply with QP (Container: 89, 50 and 10\,kbit at QP 18, 22 and 28) and grows
with resolution (City 4CIF: about 110\,kbit). The scan exposed a defect in the
carrier analysis at non-CIF resolutions, which mapped macroblocks to frames
with the CIF size and silently discarded three quarters of each 4CIF frame; it
failed closed and is fixed.
\begin{table}[t]
\centering
\caption{Stable, patchability-checked carrier capacity.}
\label{tab:capacity}
\resizebox{\linewidth}{!}{%
\begin{tabular}{lllllll}
\toprule
Video & Frames & Min bits/100 fr. & Mean bits/100 fr. & Total bits & Registered fits & In-band fits \\
\midrule
city\_4cif\_q18\_g1.h264 & 600 & 208761 & 248063 & 744190 & yes & yes \\
city\_4cif\_q22\_g1.h264 & 600 & 68084 & 82265 & 246794 & yes & yes \\
city\_4cif\_q28\_g1.h264 & 600 & 4405 & 5387 & 16162 & no & no \\
city\_cif\_q22\_g1.h264 & 300 & 537034 & 537034 & 537034 & yes & yes \\
city\_cif\_q22\_g1.h264 & 300 & 21697 & 23836 & 71508 & yes & no \\
coastguard\_cif\_q22\_g1.h264 & 300 & 549737 & 549737 & 549737 & yes & yes \\
coastguard\_cif\_q22\_g1.h264 & 300 & 31404 & 36331 & 108994 & yes & no \\
container\_cif\_q18\_g1.h264 & 300 & 86320 & 88305 & 264914 & yes & yes \\
container\_cif\_q22\_g1.h264 & 300 & 467583 & 467583 & 467583 & yes & yes \\
container\_cif\_q18\_g1.h264 & 300 & 86320 & 88305 & 264914 & yes & yes \\
container\_cif\_q22\_g1.h264 & 300 & 47866 & 49768 & 149305 & yes & yes \\
container\_cif\_q28\_g1.h264 & 300 & 9110 & 10015 & 30044 & no & no \\
crew\_4cif\_q18\_g1.h264 & 600 & 19605 & 34710 & 208259 & yes & yes \\
crew\_4cif\_q22\_g1.h264 & 600 & 11093 & 13236 & 39707 & no & no \\
crew\_4cif\_q22\_g1.h264 & 600 & 5613 & 10626 & 63756 & no & no \\
football\_cif\_q22\_g1.h264 & 260 & 14295 & 21554 & 64663 & no & no \\
hall\_monitor\_cif\_q22\_g1.h264 & 300 & 468519 & 468519 & 468519 & yes & yes \\
hall\_monitor\_cif\_q22\_g1.h264 & 300 & 27065 & 27611 & 82832 & yes & no \\
harbour\_4cif\_q22\_g1.h264 & 600 & 48816 & 49976 & 199905 & yes & yes \\
harbour\_4cif\_q22\_g1.h264 & 600 & 48816 & 49357 & 148072 & yes & yes \\
akiyo\_cif\_q22\_g1.h264 & 300 & 10501 & 10501 & 10501 & no & no \\
deadline\_cif\_q22\_g1.h264 & 1374 & 46589 & 46589 & 46589 & no & no \\
foreman\_cif\_q22\_g1.h264 & 300 & 14071 & 14071 & 14071 & no & no \\
soccer\_4cif\_q22\_g1.h264 & 600 & 3501 & 30010 & 90030 & yes & no \\
soccer\_4cif\_q22\_g1.h264 & 600 & 38552 & 43264 & 86529 & yes & no \\
\bottomrule
\end{tabular}}
\end{table}


\subsection{End-to-end embedding and verification}
Every run in Table~\ref{tab:e2e} verifies from the stego video alone. On the
modified frames the luma PSNR against the cover is 40--45\,dB, with a worst
frame of 32--39\,dB; whole-video SSIM stays above 0.999. Embedding takes
6--15 minutes and verification 2--10 minutes per clip, dominated by the Python
CAVLC analysis; proving and verifying take about 0.12\,s and 0.02\,s. The
in-band mode was also exercised end to end with a linkage proof between the
commitment extracted from the video and a separately committed credential,
and with a later disclosure of the hidden payload.
\begin{table}[t]
\centering
\caption{End-to-end embedding and video-only verification.}
\label{tab:e2e}
\resizebox{\linewidth}{!}{%
\begin{tabular}{lllllllllllll}
\toprule
Case & Mode & Envelope (B) & Carriers & Verified & Replay rej. & Tamper rej. & PSNR mod. (dB) & Worst (dB) & Frames mod. & SSIM & Embed (s) & Verify (s) \\
\midrule
city4cif\_q18\_reg & registered & 7589 & 61008 & yes & -- & -- & 47.67 & 42.82 & 100/600 & 0.99998 & 160 & 115 \\
city4cif\_q22\_reg & registered & 7589 & 61008 & yes & yes & yes & 43.66 & 39.31 & 100/600 & 0.99995 & 881 & 584 \\
city4cif\_q22\_reg\_d25 & registered & 7589 & 61896 & yes & -- & -- & 49.32 & 38.17 & 400/600 & 0.99994 & 3248 & 1309 \\
city4cif\_q22\_reg\_d50 & registered & 7589 & 61304 & yes & -- & -- & 45.36 & 38.61 & 200/600 & 0.99994 & 1443 & 698 \\
city\_cif\_q22\_reg & registered & 7589 & 61600 & yes & -- & -- & 45.35 & 38.81 & 300/300 & 0.99969 & 903 & 286 \\
city\_cif\_q22\_reg\_t1 & registered & 7589 & 61008 & yes & -- & -- & 44.09 & 40.65 & 100/300 & 0.99937 & 1338 & 816 \\
city\_cif\_q22\_reg\_t1\_spread & registered & 7589 & 61600 & yes & -- & -- & 48.79 & 44.57 & 300/300 & 0.99937 & 2197 & 2395 \\
coastguard3000f\_q22\_inband\_link & in-band & 14309 & 118024 & yes & yes & yes & 45.10 & 36.02 & 1200/3000 & 0.99992 & 5409 & 1707 \\
coastguard3000f\_q22\_reg & registered & 7589 & 62784 & yes & yes & yes & 45.26 & 35.66 & 694/3000 & 0.99996 & 2164 & 740 \\
coastguard\_cif\_q22\_reg & registered & 7589 & 61304 & yes & -- & -- & 39.97 & 33.48 & 200/300 & 0.99958 & 702 & 271 \\
coastguard\_cif\_q22\_reg\_t1 & registered & 7589 & 61008 & yes & -- & -- & 41.01 & 36.80 & 100/300 & 0.99877 & 857 & 520 \\
coastguard\_cif\_q22\_reg\_t1\_spread & registered & 7589 & 61304 & yes & -- & -- & 43.28 & 38.49 & 200/300 & 0.99805 & 1570 & 1057 \\
container\_cif\_q18\_inband & in-band & 14309 & 115064 & yes & -- & -- & 44.12 & 38.42 & 200/300 & 0.99980 & 678 & 465 \\
container\_cif\_q18\_reg & registered & 7589 & 61008 & yes & -- & -- & 43.94 & 36.41 & 100/300 & 0.99989 & 334 & 127 \\
container\_cif\_q22\_inband & in-band & 14309 & 115360 & yes & yes & yes & 40.26 & 32.13 & 300/300 & 0.99949 & 719 & 466 \\
container\_cif\_q22\_reg & registered & 7589 & 61304 & yes & yes & yes & 42.07 & 34.02 & 200/300 & 0.99975 & 580 & 182 \\
container\_cif\_q22\_reg\_t1 & registered & 7589 & 61008 & yes & -- & -- & 42.34 & 39.15 & 100/300 & 0.99855 & 892 & 558 \\
container\_cif\_q22\_reg\_t1\_spread & registered & 7589 & 61304 & yes & -- & -- & 45.69 & 41.48 & 200/300 & 0.99848 & 1275 & 788 \\
container\_cif\_q22\_reg\_t1v2\_spread & registered & 7589 & 61304 & yes & -- & -- & 45.20 & 40.21 & 200/300 & 0.99844 & 1015 & 953 \\
crew4cif\_q18\_reg & registered & 7589 & 61896 & yes & -- & -- & 48.81 & 36.65 & 400/600 & 0.99990 & 478 & 216 \\
deadline\_cif\_q22\_reg & registered & 7589 & 62784 & yes & yes & yes & 44.99 & 34.19 & 684/1000 & 0.99984 & 1793 & 1261 \\
hall\_cif\_q22\_reg & registered & 7589 & 61600 & yes & -- & -- & 41.48 & 32.72 & 300/300 & 0.99911 & 613 & 238 \\
hall\_cif\_q22\_reg\_t1\_spread & registered & 7589 & 61600 & yes & -- & -- & 47.25 & 41.01 & 300/300 & 0.99749 & 2329 & 1398 \\
hall\_monitor\_cif\_q22\_reg\_t1 & registered & 7589 & 61008 & yes & -- & -- & 42.25 & 37.74 & 100/300 & 0.99784 & 887 & 519 \\
harbour4cif\_q22\_inband & in-band & 14309 & 115360 & yes & -- & -- & 41.16 & 33.15 & 300/600 & 0.99987 & 3760 & 1108 \\
harbour4cif\_q22\_reg & registered & 7589 & 61304 & yes & -- & -- & 42.83 & 32.91 & 200/600 & 0.99994 & 2357 & 1954 \\
harbour4cif\_q22\_reg\_d25 & registered & 7589 & 62488 & yes & -- & -- & 50.74 & 38.47 & 595/600 & 0.99995 & 915 & 414 \\
harbour4cif\_q22\_reg\_d50 & registered & 7589 & 61600 & yes & -- & -- & 45.95 & 34.98 & 300/600 & 0.99994 & 471 & 215 \\
soccer4cif\_q22\_reg\_s2 & registered & 7589 & 61304 & yes & -- & -- & 44.64 & 37.11 & 200/600 & 0.99991 & 310 & 142 \\
\bottomrule
\end{tabular}}
\end{table}


\subsection{Steganalysis}
Classical targeted tests against the exact cover (Westfeld--Pfitzmann
pairs-of-values on the carrier magnitudes, a two-sample magnitude histogram
test, trailing-one sign tests and SPA) do not separate cover and stego
(Table~\ref{tab:steg-classical}). A learned detector is more effective. We
train Fisher-linear-discriminant ensembles~\cite{ref:ensemble} on 340
coefficient-domain features (per-position magnitude histograms, parity of
large magnitudes, adjacent magnitude co-occurrences and sign shares) and on
500 reduced spatial-rich-model features~\cite{ref:srm}, using only frames the
embedding touched. With leave-one-clip-out evaluation over five CIF clips
(1\,300 frame pairs) the detection error is $P_E=0.40$ for coefficient
features and $0.45$ for pixel features; an attacker trained on the same clip
reaches $P_E=0.36$ (Table~\ref{tab:steg-ml}). These are per-frame figures;
aggregating evidence over all modified frames of a video would detect the
embedding more reliably. Spreading the envelope over more frames lowers the
per-frame density and raises $P_E$ (Section~\ref{sec:density}).
\begin{table}[t]
\centering
\caption{Targeted classical steganalysis against the exact cover.}
\label{tab:steg-classical}
\resizebox{\linewidth}{!}{%
\begin{tabular}{lllllll}
\toprule
Case & PoV P(embed) cover & PoV P(embed) stego & Magnitude hist. p & T1 sign p & SPA cover & SPA stego \\
\midrule
city4cif\_q18\_reg & 0.000 & 0.000 & 1.000 & 1.000 & 0.41693 & 0.41697 \\
city4cif\_q22\_reg & 0.000 & 0.000 & 0.998 & 1.000 & 0.39448 & 0.39456 \\
city4cif\_q22\_reg\_d25 & 0.000 & 0.000 & 1.000 & 1.000 & 0.39448 & 0.39450 \\
city4cif\_q22\_reg\_d50 & 0.000 & 0.000 & 1.000 & 1.000 & 0.39448 & 0.39451 \\
city\_cif\_q22\_reg & 0.000 & 0.000 & 1.000 & 1.000 & 0.40765 & 0.40772 \\
city\_cif\_q22\_reg\_t1 & 0.000 & 0.000 & 1.000 & 0.119 & 0.40765 & 0.40857 \\
city\_cif\_q22\_reg\_t1\_spread & 0.000 & 0.000 & 1.000 & 0.475 & 0.40765 & 0.40803 \\
coastguard3000f\_q22\_reg & 0.000 & 0.000 & 1.000 & 1.000 & 0.40661 & 0.40664 \\
coastguard\_cif\_q22\_reg & 0.000 & 0.000 & 1.000 & 1.000 & 0.40661 & 0.40671 \\
coastguard\_cif\_q22\_reg\_t1 & 0.000 & 0.000 & 1.000 & 0.176 & 0.40661 & 0.40704 \\
coastguard\_cif\_q22\_reg\_t1\_spread & 0.000 & 0.000 & 1.000 & 0.805 & 0.40661 & 0.40684 \\
container\_cif\_q18\_inband & 0.000 & 0.000 & 1.000 & 1.000 & 0.33144 & 0.33156 \\
container\_cif\_q18\_reg & 0.000 & 0.000 & 1.000 & 1.000 & 0.33144 & 0.33165 \\
container\_cif\_q22\_inband & 0.000 & 0.000 & 1.000 & 1.000 & 0.27711 & 0.27752 \\
container\_cif\_q22\_reg & 0.000 & 0.000 & 1.000 & 1.000 & 0.27711 & 0.27729 \\
container\_cif\_q22\_reg\_t1 & 0.000 & 0.000 & 1.000 & 0.061 & 0.27711 & 0.27788 \\
container\_cif\_q22\_reg\_t1\_spread & 0.000 & 0.000 & 1.000 & 0.283 & 0.27711 & 0.27759 \\
container\_cif\_q22\_reg\_t1v2\_spread & 0.000 & 0.000 & 1.000 & 0.139 & 0.27711 & 0.27753 \\
crew4cif\_q18\_reg & 0.000 & 0.000 & 1.000 & 1.000 & 0.31194 & 0.31198 \\
deadline\_cif\_q22\_reg & 0.000 & 0.000 & 1.000 & 1.000 & 0.28793 & 0.28810 \\
hall\_cif\_q22\_reg & 0.000 & 0.000 & 1.000 & 1.000 & 0.26447 & 0.26538 \\
hall\_cif\_q22\_reg\_t1\_spread & 0.000 & 0.000 & 1.000 & 0.419 & 0.25033 & 0.25206 \\
hall\_monitor\_cif\_q22\_reg\_t1 & 0.000 & 0.000 & 1.000 & 0.173 & 0.25033 & 0.25419 \\
harbour4cif\_q22\_inband & 0.000 & 0.000 & 1.000 & 1.000 & 0.42423 & 0.42431 \\
harbour4cif\_q22\_reg & 0.000 & 0.000 & 1.000 & 1.000 & 0.42423 & 0.42435 \\
harbour4cif\_q22\_reg\_d25 & 0.000 & 0.000 & 1.000 & 1.000 & 0.42423 & 0.42422 \\
harbour4cif\_q22\_reg\_d50 & 0.000 & 0.000 & 1.000 & 1.000 & 0.42423 & 0.42425 \\
soccer4cif\_q22\_reg\_s2 & 0.000 & 0.000 & 0.997 & 1.000 & 0.30408 & 0.30414 \\
\bottomrule
\end{tabular}}
\end{table}

\begin{table}[t]
\centering
\caption{Learned steganalysis on modified frames (P\_E = 0.5 is chance).}
\label{tab:steg-ml}
\resizebox{\linewidth}{!}{%
\begin{tabular}{llll}
\toprule
Protocol & Features / detector & P\_E & AUC \\
\midrule
paired\_half\_split & coefficient/logistic & 0.414 & 0.617 \\
paired\_half\_split & coefficient/fld\_ensemble & 0.394 & 0.638 \\
paired\_half\_split & pixel/logistic & 0.424 & 0.608 \\
paired\_half\_split & pixel/fld\_ensemble & 0.410 & 0.627 \\
paired\_half\_split & combined/logistic & 0.401 & 0.633 \\
paired\_half\_split & combined/fld\_ensemble & 0.357 & 0.685 \\
leave\_one\_clip\_out & coefficient/logistic & 0.407 & 0.604 \\
leave\_one\_clip\_out & coefficient/fld\_ensemble & 0.399 & 0.628 \\
leave\_one\_clip\_out & pixel/logistic & 0.454 & 0.538 \\
leave\_one\_clip\_out & pixel/fld\_ensemble & 0.458 & 0.536 \\
leave\_one\_clip\_out & combined/logistic & 0.437 & 0.575 \\
leave\_one\_clip\_out & combined/fld\_ensemble & 0.422 & 0.583 \\
\bottomrule
\end{tabular}}
\end{table}


\subsection{Carrier choice}
\label{sec:carrier}
We compared the magnitude-parity carrier with the trailing-one (T1) sign flip
used by earlier CAVLC schemes, through the same pipeline, on four CIF clips,
both at the same frame count as the magnitude carrier and packed into 100
frames. All twelve runs verify. Sign carriers are an order of magnitude more
plentiful (470--550\,kbit per 100 frames) and always patch as a single bit.
At equal frame count they give higher luma PSNR on modified frames
(43.3--48.8 vs.\ 40.0--45.4\,dB; worst frame 38.5--44.6 vs.\ 32.7--38.8\,dB) but
slightly lower whole-video SSIM, and they are much harder to see in the
coefficient domain (leave-one-clip-out $P_E=0.47$ vs.\ $0.39$) yet easier in
the pixel domain ($0.40$ vs.\ $0.46$); combined features give $0.40$ vs.\
$0.42$. Packed into 100 frames, sign flips are clearly detectable ($P_E=0.31$
with pixel features, and a per-position sign test flags one clip at
$p=0.001$). Verification with sign carriers is 4--8 times slower because the
reduced stable-carrier analysis does not apply. We keep the magnitude carrier
as the default for its verification speed and pixel-domain behaviour, and
expose the sign carrier as an option when quality matters more than speed.

\subsection{Density}
\label{sec:density}
Simulated on the same clips and detector, using 50, 25 and 10\,\% of the
carriers of each segment raises the coefficient-feature $P_E$ from 0.32 to
0.40, 0.44 and 0.47. LSB matching instead of LSB replacement is more
detectable ($P_E=0.26$) and was not adopted. Measured density results on the
4CIF clips are reported in the artifact.

\subsection{Comparison with baselines}
Table~\ref{tab:baselines} compares proof systems that could carry the claim.
Groth16 is two orders of magnitude smaller but not post-quantum and needs a
circuit-specific trusted setup. ML-DSA-44 is post-quantum and smaller but
proves nothing about a hidden payload. ISW21 is a lattice zkSNARK with a
designated verifier whose knowledge assumption is broken by quantum oblivious
sampling.
\begin{table}[t]
\centering
\caption{Comparison with baseline proof and signature schemes.}
\label{tab:baselines}
\resizebox{\linewidth}{!}{%
\begin{tabular}{lllllll}
\toprule
Scheme & PQ & ZK & In-video bytes & Prove (ms) & Verify (ms) & Setup \\
\midrule
BDLOP (this work, registered) & yes & yes & 7589 & 187 & 41.9 & transparent \\
BDLOP (this work, in-band) & yes & yes & 14309 & 187 & 41.9 & transparent \\
Groth16 (BN128) & no & yes & 147 & 7222 & 2550 & trusted, per circuit \\
ML-DSA-44 (signature) & yes & no & 2420 & 18.6 & 0.33 & none \\
ISW21 (recorded) & broken KS assumption & stat. 2^-40 & 23907 & 328 & 0.6 & trusted, designated verifier \\
\bottomrule
\end{tabular}}
\end{table}


\section{Discussion and limitations}
\label{sec:discussion}

Soundness is relaxed; payload binding holds with respect to relaxed openings.
The implementation is not constant time and has not been audited. The
carrier order is keyed by a symmetric synchronization (stego) key shared with
the verifier. This is inherent to steganography: without a secret, anyone who
knows the scheme can locate the carriers and test for an envelope. The key
affects only where bits are found; soundness and zero knowledge of the proof
do not depend on it. Where undetectability is not required, the key can be
derived from the public carrier profile, making verification public. Proofs that each payload bit is
binary require a fully splitting ring and the product proofs of
ALS20/ENS20. Capacity limits the scheme to clips of a few hundred frames at
CIF and fewer at higher resolution.

\section{Conclusion}

We showed that a post-quantum zero-knowledge proof can travel inside an
H.264 bitstream, bound to that bitstream and to a verifier session, and be
verified from the video alone.

\bibliographystyle{IEEEtran}
\bibliography{references}

\end{document}
